% Secao 2: O que e Redes de Computadores? (A magica por tras da internet)

% Slide 4: Como a Internet Funciona?
\begin{frame}{Você Já Pensou em Como a Internet Funciona?}{\faWifi\ Da tela do seu celular para qualquer lugar do planeta em segundos}
  \begin{columns}[t]
    \begin{column}{0.48\textwidth}
      \begin{bluecard}{\faMobile\ O Que Você Vê Todo Santo Dia}
        \scriptsize
        \begin{itemize}\setlength{\itemsep}{0.18em}
          \item \textbf{Games Online sem Lag:} Como Free Fire e Roblox sincronizam dezenas de jogadores ao vivo?
          \item \textbf{Vídeos em 4K:} Como streams no YouTube e TikTok chegam sem travar no seu celular?
          \item \textbf{Mensagens Instantâneas:} Como uma foto no WhatsApp atravessa o mundo em milissegundos?
          \item \textbf{Segurança no PIX:} Como seu dinheiro viaja com criptografia e proteção antifraude?
        \end{itemize}
      \end{bluecard}
    \end{column}
    
    \begin{column}{0.48\textwidth}
      \begin{greencard}{\faProjectDiagram\ O Que Existe nos Bastidores}
        \scriptsize
        \begin{itemize}\setlength{\itemsep}{0.18em}
          \item \textbf{Espinha Dorsal Global:} Milhares de km de fibra óptica, cabos submarinos e satélites.
          \item \textbf{Roteadores Inteligentes:} Equipamentos que guiam cada bit na velocidade da luz.
          \item \textbf{Data Centers Gigantescos:} Servidores 24h por dia sustentando a Nuvem e a IA.
          \item \textbf{Sem Redes, Nada Disso Existe:} O celular vira só uma calculadora sem a rede de dados!
        \end{itemize}
      \end{greencard}
    \end{column}
  \end{columns}
  
  \vspace{0.1cm}
  \begin{center}
    \scriptsize \textcolor{IFCEDarkGreen}{\textbf{\faGlobe\ Redes de Computadores é a infraestrutura invisível que conecta a humanidade!}}
  \end{center}
\end{frame}

% Slide 5: Os 4 Superpoderes do Tecnico em Redes
\begin{frame}{O Que Faz um Técnico em Redes?}{\faUserAstronaut\ Os 4 pilares práticos da profissão que conecta o mundo}
  \begin{columns}[t]
    \begin{column}{0.24\textwidth}
      \begin{bluecard}{\faNetworkWired\ 1. Conexão}
        \tiny
        \textbf{Infraestrutura}
        \begin{itemize}\setlength{\itemsep}{0.12em}
          \item Redes Wi-Fi velozes.
          \item Cabeamento e fibra.
          \item Interliga empresas.
        \end{itemize}
      \end{bluecard}
    \end{column}
    
    \begin{column}{0.24\textwidth}
      \begin{greencard}{\faServer\ 2. Servidores}
        \tiny
        \textbf{Coração da Nuvem}
        \begin{itemize}\setlength{\itemsep}{0.12em}
          \item Linux e Windows.
          \item Gestão de usuários.
          \item Mantém serviços no ar.
        \end{itemize}
      \end{greencard}
    \end{column}
    
    \begin{column}{0.24\textwidth}
      \begin{redcard}{\faLock\ 3. Segurança}
        \tiny
        \textbf{Defesa Cibernética}
        \begin{itemize}\setlength{\itemsep}{0.12em}
          \item Bloqueia malwares.
          \item Configura firewalls.
          \item Protege senhas e dados.
        \end{itemize}
      \end{redcard}
    \end{column}
    
    \begin{column}{0.24\textwidth}
      \begin{purplecard}{\faCode\ 4. Código}
        \tiny
        \textbf{Automação \& Web}
        \begin{itemize}\setlength{\itemsep}{0.12em}
          \item Scripts em Python.
          \item Páginas Web e bancos.
          \item Soluções IoT e sensores.
        \end{itemize}
      \end{purplecard}
    \end{column}
  \end{columns}
  
  \vspace{0.1cm}
  \begin{statcard}
    \tiny \textcolor{TechDark}{\faWrench\ \textbf{Profissão Completa:} Você aprende tanto a parte física (hardware, cabos, placas) quanto a parte lógica (sistemas, códigos, segurança)!}
  \end{statcard}
\end{frame}
